\bta{2025}

\noindent
{\kaishu\color{gray!50!black}
（本年试题为考生回忆版，题面文字与个别记号可能与原卷略有出入，
凡有疑义处已在解答中注明。）}

\vspace*{2ex}

\begin{enumerate}
	%\renewcommand{\labelenumi}{\arabic{enumi}.}
	% A(\Alph) a(\alph) I(\Roman) i(\roman) 1(\arabic)
	%设定全局标号series=example	%引用全局变量resume=example
	%[topsep=-0.3em,parsep=-0.3em,itemsep=-0.3em,partopsep=-0.3em]
	%可使用leftmargin调整列表环境左边的空白长度 [leftmargin=0em]
	\item
选择题
\begin{enumerate}
	%\renewcommand{\labelenumi}{\arabic{enumi}.}
	% A(\Alph) a(\alph) I(\Roman) i(\roman) 1(\arabic)
	%设定全局标号series=example	%引用全局变量resume=example
	%[topsep=-0.3em,parsep=-0.3em,itemsep=-0.3em,partopsep=-0.3em]
	%可使用leftmargin调整列表环境左边的空白长度 [leftmargin=0em]
	\item
求极限 $\lim\limits_{n \rightarrow \infty} \frac{\ln \left(1+\frac{1}{2}+\cdots+\frac{1}{n}\right)}{n}$ \xzanswer{D}

\fourchoices
{$1$}
{$-1$}
{$e$}
{$0$}


\item
下列何者为最低阶无穷小（$x \rightarrow 0$） \xzanswer{C}

\fourchoices
{$1-\frac{\sin x}{x}$}
{$x \tan x$}
{$\ln (1+x)$}
{$\frac{x\left(2^x-1\right)}{x-1}$}


\item
设 $x f^{\prime \prime}(x)+3 f^{\prime}(x)=x-x e^{-x}$，存在 $x_0$ 使得 $f^{\prime}\left(x_0\right)=0$，
则该点是 \xzanswer{C}

\fourchoices
{极大值点}
{极小值点}
{拐点}
{都不是}


\item
设 $F(x, y, z)=0$，其中任意一个变量都可以改写成另外两个变量的函数，
记号 $\left(\frac{d y}{d x}\right)_z$ 表示固定 $z$ 对 $x$ 求导。设各偏导数在某点均不为零，求
\[
\left(\frac{d y}{d x}\right)_z \cdot\left(\frac{d x}{d z}\right)_y \cdot\left(\frac{d z}{d y}\right)_x
\]
\xzanswer{A}

\fourchoices
{$-1$}
{$1$}
{$0$}
{$x y z$}


\item
点 $A(1,0,1)$，$B(0,1,1)$，$C(a, 0,-1)$ 构成三角形 $ABC$，
求 $a$ 为何值时 $S_{\triangle A B C}$ 取最小值，并求 $S_{\min }$ 的值 \xzanswer{C}
（原卷各选项后另留有填写 $S_{\min }$ 的横线，答案为 \tk{$S_{\min }=\sqrt{2}$}）

\fourchoices
{$a=0$}
{$a=\sqrt{2}$}
{$a=1$}
{$a=2$}


\item
设 $f(x)$ 满足 $f(x+1)=a f(x)$ 恒成立，且在 $x=0$ 处可导，$f^{\prime}(0)=b$，
则 $x=1$ 处的导数情况为 \xzanswer{D}

\fourchoices
{不存在}
{$a$}
{$b$}
{$a b$}


\item
将下列积分按 $x \rightarrow 0^{+}$ 时由低阶到高阶无穷小排序：
\[
a=\int_0^x \cos t^5 d t,\qquad b=\int_0^{x^2} \ln (1+\sqrt{t}) d t,\qquad
c=\int_0^{\sqrt{x}} t \sin t^2 d t .
\]
答： \tk{$a$，$c$，$b$}


\item
微分方程 $y^{\prime}+P(x) y=Q(x)$ 的两个特解为 $y_1(x)$，$y_2(x)$，其通解为

\hspace*{2em}\tk{$y_1(x)+C\left(y_1(x)-y_2(x)\right)$}


\item
幂级数 $\sum\limits_{n=0}^{\infty} \sqrt{n!}\, x^n$ 的收敛半径 \xzanswer{B}

\threechoices
{$R=1$}
{$R=0$}
{$R=+\infty$}


\item
下列积分发散的是 \xzanswer{D}

\fourchoices
{$\int_1^{+\infty} \frac{\ln x}{(1+x)^{3 / 2}} d x$}
{$\int_0^{+\infty} \frac{\ln x}{1-x^2} d x$}
{$\int_1^{+\infty} \frac{\ln ^2 x}{x^2} d x$}
{$\int_0^{+\infty} \frac{1}{\sqrt{\sin x}} d x$}


\end{enumerate}


\banswer{
\begin{enumerate}
	%\renewcommand{\labelenumi}{\arabic{enumi}.}
	% A(\Alph) a(\alph) I(\Roman) i(\roman) 1(\arabic)
	%设定全局标号series=example	%引用全局变量resume=example
	%[topsep=-0.3em,parsep=-0.3em,itemsep=-0.3em,partopsep=-0.3em]
	%可使用leftmargin调整列表环境左边的空白长度 [leftmargin=0em]
	\item
	\textbf{D}。记 $H_n=1+\frac{1}{2}+\cdots+\frac{1}{n}$，则 $H_n \sim \ln n$，
	故 $\ln H_n \sim \ln \ln n$，而 $\frac{\ln \ln n}{n} \rightarrow 0$。

	\item
	\textbf{C}。逐个定阶：
	$1-\frac{\sin x}{x}=1-\left(1-\frac{x^2}{6}+o\left(x^2\right)\right) \sim \frac{x^2}{6}$，二阶；
	$x \tan x \sim x^2$，二阶；
	$\ln (1+x) \sim x$，\textbf{一阶}；
	$\frac{x\left(2^x-1\right)}{x-1} \sim \frac{x \cdot x \ln 2}{-1}=-x^2 \ln 2$，二阶。
	阶数最低者为 C。

	\item
	\textbf{C}。方程两端乘 $x^2$：$x^3 f^{\prime \prime}+3 x^2 f^{\prime}=x^3\left(1-e^{-x}\right)$，即
	\[
	\left(x^3 f^{\prime}(x)\right)^{\prime}=x^3\left(1-e^{-x}\right)=: g(x).
	\]
	因 $x$ 与 $1-e^{-x}$ 同号，故 $g(x)>0\ (x \neq 0)$，$g(0)=0$。
	令 $H(x)=x^3 f^{\prime}(x)$，由方程在 $x=0$ 处取值得 $3 f^{\prime}(0)=0$，故 $H(0)=0$，
	且 $H^{\prime}=g \geqslant 0$ 且仅在一点为零，故 $H$ 严格递增。
	于是 $x>0$ 时 $H>0$，$x<0$ 时 $H<0$，从而
	\[
	f^{\prime}(x)=\frac{H(x)}{x^3}>0,\qquad \forall x \neq 0 .
	\]
	可见 $f^{\prime}$ 只能在 $x=0$ 处为零，即 $x_0=0$。
	又 $1-e^{-t}=t-\frac{t^2}{2}+o\left(t^2\right)$ 给出
	$H(x)=\int_0^x t^3\left(1-e^{-t}\right) d t=\frac{x^5}{5}-\frac{x^6}{12}+o\left(x^6\right)$，故
	\[
	f^{\prime}(x)=\frac{x^2}{5}-\frac{x^3}{12}+o\left(x^3\right),\qquad
	f^{\prime \prime}(x)=\frac{2 x}{5}-\frac{x^2}{4}+o\left(x^2\right).
	\]
	因 $f^{\prime}>0$（$x \neq 0$），$x_0=0$ 不是极值点；
	而 $f^{\prime \prime}$ 在 $0$ 处变号（左负右正），故 $x_0=0$ 是\textbf{拐点}。

	\item
	\textbf{A}。由隐函数求导法则
	$\left(\frac{d y}{d x}\right)_z=-\frac{F_x}{F_y}$，
	$\left(\frac{d x}{d z}\right)_y=-\frac{F_z}{F_x}$，
	$\left(\frac{d z}{d y}\right)_x=-\frac{F_y}{F_z}$，三者相乘得
	\[
	(-1)^3 \cdot \frac{F_x}{F_y} \cdot \frac{F_z}{F_x} \cdot \frac{F_y}{F_z}=-1 .
	\]

	\item
	\textbf{C}。$\overrightarrow{A B}=(-1,1,0)$，$\overrightarrow{A C}=(a-1,0,-2)$，
	\[
	\overrightarrow{A B} \times \overrightarrow{A C}=(-2,\ -2,\ 1-a),\qquad
	\left|\overrightarrow{A B} \times \overrightarrow{A C}\right|^2=8+(a-1)^2 .
	\]
	故 $S=\frac{1}{2} \sqrt{8+(a-1)^2}$ 在 $a=1$ 时取最小值
	$S_{\min }=\frac{1}{2} \sqrt{8}=\sqrt{2}$。

	\item
	\textbf{D}。$f(1+h)=a f(h)$，$f(1)=a f(0)$，故
	\[
	\lim _{h \rightarrow 0} \frac{f(1+h)-f(1)}{h}=a \lim _{h \rightarrow 0} \frac{f(h)-f(0)}{h}=a f^{\prime}(0)=a b .
	\]

	\item
	\textbf{$a$，$c$，$b$}。逐个定阶（$x \rightarrow 0^{+}$）：
	\[
	a=\int_0^x \cos t^5 d t \sim x \quad(\text{一阶}),
	\]
	\[
	b=\int_0^{x^2} \ln (1+\sqrt{t}) d t \sim \int_0^{x^2} \sqrt{t}\, d t=\frac{2}{3} x^3 \quad(\text{三阶}),
	\]
	\[
	c=\int_0^{\sqrt{x}} t \sin t^2 d t=\frac{1}{2} \int_0^x \sin u\, d u=\frac{1-\cos x}{2} \sim \frac{x^2}{4} \quad(\text{二阶}).
	\]
	故由低阶到高阶为 $a$，$c$，$b$。

	\item
	$y_1(x)+C\left(y_1(x)-y_2(x)\right)$。这里需 $y_1\not\equiv y_2$。此时 $y_1-y_2$ 是齐次方程 $y^{\prime}+P(x) y=0$ 的非零解，
	$y_1$ 是原非齐次方程的一个特解，由“非齐次特解 $+$ 齐次通解”即得。

	\item
	\textbf{B}。由 $\lim\limits_{n \rightarrow \infty} \frac{\sqrt{(n+1)!}}{\sqrt{n!}}=\lim\limits_{n \rightarrow \infty} \sqrt{n+1}=+\infty$
	知 $\frac{1}{R}=+\infty$，即 $R=0$（级数仅在 $x=0$ 处收敛）。

	\item
	\textbf{D}。A 中 $x \rightarrow+\infty$ 时被积函数 $\sim \frac{\ln x}{x^{3 / 2}}$，收敛；
	C 中 $\int_1^{+\infty} \frac{\ln ^2 x}{x^2} d x=2$，收敛；
	B 中 $x \rightarrow 1$ 时 $\ln x \sim x-1$，$\frac{\ln x}{1-x^2} \rightarrow-\frac{1}{2}$ 为可去奇性，
	$x \rightarrow 0^{+}$ 时 $\ln x$ 可积，$x \rightarrow+\infty$ 时 $\sim-\frac{\ln x}{x^2}$，故收敛。
	D 中 $\sin x$ 在 $x>\pi$ 时变号（被积函数已无意义），
	即便只在各 $\sin x>0$ 的区间上取值，每个 $[2k\pi,(2k+1)\pi]$（$k=0,1,\ldots$） 上的积分为同一正常数
	$\int_0^\pi \frac{d x}{\sqrt{\sin x}}>0$，无穷多个相加必发散。
\end{enumerate}
}


\item
已知
$f(x, y)=\left\{\begin{array}{cl}
\frac{\sin ^2 x \sin y+\sin ^2 y \sin x}{x^2-x y+y^2}, & x^2+y^2 \neq 0 \\
1, & x^2+y^2=0
\end{array}\right.$，
判断（讨论）$f(x, y)$ 在点 $(0,0)$ 处的连续性。

\banswer{
	\textbf{分母的下界.} 对一切 $(x, y) \neq(0,0)$，
	\[
	x^2-x y+y^2=\frac{1}{2}\left(x^2+y^2\right)+\frac{1}{2}(x-y)^2 \geqslant \frac{1}{2}\left(x^2+y^2\right)=\frac{\rho^2}{2}>0,
	\]
	其中 $\rho=\sqrt{x^2+y^2}$。特别地分母恒不为零，$f$ 在原点外处处有定义。

	\textbf{分子的上界.} 由 $|\sin u| \leqslant|u|$ 及 $|x|,|y| \leqslant \rho$，
	\[
	\left|\sin ^2 x \sin y+\sin ^2 y \sin x\right| \leqslant x^2|y|+y^2|x| \leqslant 2 \rho^3 .
	\]

	\textbf{极限.} 于是
	\[
	|f(x, y)| \leqslant \frac{2 \rho^3}{\rho^2 / 2}=4 \rho \longrightarrow 0 \quad(\rho \rightarrow 0),
	\]
	故 $\lim\limits_{(x, y) \rightarrow(0,0)} f(x, y)=0$。

	\textbf{结论.} 而 $f(0,0)=1 \neq 0$，故 $f$ 在 $(0,0)$ 处\textbf{不连续}；
	该间断点是可去间断点——若把定义改为 $f(0,0)=0$，则 $f$ 在 $(0,0)$ 处连续。
}


\item
设 $f(x)$ 在 $[0,2]$ 上二阶可导，$f(0)=f(2)=0$，$F(x)=(x-1)^2 f(x)$，
证明至少存在一点 $\xi \in(0,2)$，使得 $F^{\prime \prime}(\xi)=0$。

\banswer{
	\textbf{三个零点.} 由 $F(x)=(x-1)^2 f(x)$ 得
	\[
	F(0)=1 \cdot f(0)=0,\qquad F(1)=0 \cdot f(1)=0,\qquad F(2)=1 \cdot f(2)=0 .
	\]
	又 $f$ 在 $[0,2]$ 上二阶可导，故 $F$ 在 $[0,2]$ 上二阶可导。

	\textbf{两次 Rolle.} 在 $[0,1]$ 上对 $F$ 用 Rolle 定理，存在 $\xi_1 \in(0,1)$ 使 $F^{\prime}\left(\xi_1\right)=0$；
	在 $[1,2]$ 上同理存在 $\xi_2 \in(1,2)$ 使 $F^{\prime}\left(\xi_2\right)=0$。

	\textbf{第三次 Rolle.} $F^{\prime}$ 在 $\left[\xi_1, \xi_2\right]$ 上连续、在 $\left(\xi_1, \xi_2\right)$ 内可导，
	且 $F^{\prime}\left(\xi_1\right)=F^{\prime}\left(\xi_2\right)=0$，再用 Rolle 定理，存在
	\[
	\xi \in\left(\xi_1, \xi_2\right) \subset(0,2),\qquad F^{\prime \prime}(\xi)=0 . \qquad \blacksquare
	\]
}


\item
计算 $I=\int_2^4 \frac{1}{\sqrt{(x-1)(5-x)}} d x$。

\banswer{
	\par\vspace{1mm}\noindent
	\textbf{注：}{\kaishu
	回忆版题面写作 $\sqrt{(1-x)(5-x)}$，但当 $x \in[2,4]$ 时 $(1-x)(5-x)<0$，
	根式无意义，应为 $\sqrt{(x-1)(5-x)}$。
	}\par

	配方：
	\[
	(x-1)(5-x)=-x^2+6 x-5=4-(x-3)^2 .
	\]
	故
	\[
	I=\int_2^4 \frac{d x}{\sqrt{4-(x-3)^2}}
	=\left[\arcsin \frac{x-3}{2}\right]_2^4
	=\arcsin \frac{1}{2}-\arcsin \left(-\frac{1}{2}\right)
	=\frac{\pi}{6}+\frac{\pi}{6}=\frac{\pi}{3} .
	\]
}


\item
证明：$I=\lim\limits_{n \rightarrow \infty}\left(1+\frac{1}{2}+\frac{1}{3}+\cdots+\frac{1}{n-1}-\ln n\right)$ 存在。
(提示 $\int_1^x \frac{1}{t} d t=\ln x$)

\banswer{
	记 $a_n=\sum\limits_{k=1}^{n-1} \frac{1}{k}-\ln n\ (n \geqslant 2)$。

	\textbf{单调递增.} 由 $\ln \frac{n+1}{n}=\int_n^{n+1} \frac{d t}{t}$ 及 $t \in(n, n+1)$ 时 $\frac{1}{t}<\frac{1}{n}$，得
	\[
	a_{n+1}-a_n=\frac{1}{n}-\ln \frac{n+1}{n}=\frac{1}{n}-\int_n^{n+1} \frac{d t}{t}>\frac{1}{n}-\frac{1}{n}=0 .
	\]
	故 $\left\{a_n\right\}$ 严格递增。

	\textbf{有上界.} 同理 $t \in(k, k+1)$ 时 $\frac{1}{t}>\frac{1}{k+1}$，故
	\[
	\ln n=\int_1^n \frac{d t}{t}=\sum_{k=1}^{n-1} \int_k^{k+1} \frac{d t}{t}>\sum_{k=1}^{n-1} \frac{1}{k+1}=\sum_{k=2}^n \frac{1}{k} .
	\]
	于是
	\[
	a_n=\sum_{k=1}^{n-1} \frac{1}{k}-\ln n<\sum_{k=1}^{n-1} \frac{1}{k}-\sum_{k=2}^n \frac{1}{k}=1-\frac{1}{n}<1 .
	\]

	\textbf{结论.} $\left\{a_n\right\}$ 单调递增且有上界 $1$，由单调有界定理，极限 $I$ 存在，
	且 $0<I \leqslant 1$。（该极限即 Euler--Mascheroni 常数 $\gamma \approx 0.5772$。）$\blacksquare$
}


\item
\begin{enumerate}
	%\renewcommand{\labelenumi}{\arabic{enumi}.}
	% A(\Alph) a(\alph) I(\Roman) i(\roman) 1(\arabic)
	%设定全局标号series=example	%引用全局变量resume=example
	%[topsep=-0.3em,parsep=-0.3em,itemsep=-0.3em,partopsep=-0.3em]
	%可使用leftmargin调整列表环境左边的空白长度 [leftmargin=0em]
	\item
	写出二元函数 $f(x, y)$ 在一点处可微的定义；
	\item
	判断
	$f(x, y)=\left\{\begin{array}{cl}
	e^{-\frac{1}{x^2+y^2}}, & x^2+y^2>0 \\
	0, & x^2+y^2=0
	\end{array}\right.$
	在 $(0,0)$ 处是否可微。
\end{enumerate}

\banswer{
\begin{enumerate}
	%\renewcommand{\labelenumi}{\arabic{enumi}.}
	% A(\Alph) a(\alph) I(\Roman) i(\roman) 1(\arabic)
	%设定全局标号series=example	%引用全局变量resume=example
	%[topsep=-0.3em,parsep=-0.3em,itemsep=-0.3em,partopsep=-0.3em]
	%可使用leftmargin调整列表环境左边的空白长度 [leftmargin=0em]
	\item
	\textbf{定义.} 设 $f$ 在点 $\left(x_0, y_0\right)$ 的某邻域内有定义。若存在常数 $A$，$B$，使得当
	$\rho=\sqrt{(\Delta x)^2+(\Delta y)^2} \rightarrow 0$ 时
	\[
	\Delta f=f\left(x_0+\Delta x, y_0+\Delta y\right)-f\left(x_0, y_0\right)=A \Delta x+B \Delta y+o(\rho),
	\]
	则称 $f$ 在 $\left(x_0, y_0\right)$ 处\textbf{可微}，
	并称 $A \Delta x+B \Delta y$ 为 $f$ 在该点的全微分；
	此时必有 $A=f_x\left(x_0, y_0\right)$，$B=f_y\left(x_0, y_0\right)$。

	\item
	\textbf{可微}，且 $\left.d f\right|_{(0,0)}=0$。

	先求两个偏导数。由定义，
	\[
	f_x(0,0)=\lim _{x \rightarrow 0} \frac{f(x, 0)-f(0,0)}{x}=\lim _{x \rightarrow 0} \frac{e^{-1 / x^2}}{x}
	\overset{s=1/|x|}{=}\lim _{s \rightarrow+\infty} \pm s e^{-s^2}=0 ,
	\]
	由对称性同理 $f_y(0,0)=0$。

	再验证增量式。取 $A=B=0$，则当 $\rho \rightarrow 0^{+}$ 时
	\[
	\frac{|\Delta f-0 \cdot \Delta x-0 \cdot \Delta y|}{\rho}=\frac{e^{-1 / \rho^2}}{\rho}
	\overset{s=1/\rho}{=} s e^{-s^2} \longrightarrow 0 \quad(s \rightarrow+\infty),
	\]
	即 $\Delta f=o(\rho)$。按定义 $f$ 在 $(0,0)$ 处可微，全微分为 $0$。

	\par\vspace{1mm}\noindent
	\textbf{注：}{\kaishu
	$e^{-1 / \rho^2}$ 在原点是“平坦”到无穷阶的：它比任何 $\rho^k$ 都小，
	因此不仅可微，事实上 $f$ 在原点还是 $C^{\infty}$ 的，各阶偏导数全为 $0$。
	}\par
\end{enumerate}
}


\item
求 $f=8 x y z$ 在条件 $\frac{x^2}{a^2}+\frac{y^2}{b^2}+\frac{z^2}{c^2}=1$（$a, b, c>0$）下的极大值。

\banswer{
	\textbf{均值不等式法.} 记 $u=\frac{x^2}{a^2}$，$v=\frac{y^2}{b^2}$，$w=\frac{z^2}{c^2}$，
	则 $u, v, w \geqslant 0$ 且 $u+v+w=1$。由
	\[
	f^2=64 x^2 y^2 z^2=64 a^2 b^2 c^2 \cdot u v w
	\]
	及三元均值不等式 $u v w \leqslant\left(\frac{u+v+w}{3}\right)^3=\frac{1}{27}$，得
	\[
	f^2 \leqslant \frac{64 a^2 b^2 c^2}{27},\qquad |f| \leqslant \frac{8 a b c}{3 \sqrt{3}}=\frac{8 \sqrt{3}}{9} a b c .
	\]
	等号当且仅当 $u=v=w=\frac{1}{3}$，即
	\[
	|x|=\frac{a}{\sqrt{3}},\qquad |y|=\frac{b}{\sqrt{3}},\qquad |z|=\frac{c}{\sqrt{3}}
	\]
	时取到；再取符号使 $x y z>0$（如全取正号），即得

	\[
	f_{\max }=\frac{8 a b c}{3 \sqrt{3}}=\frac{8 \sqrt{3}}{9} a b c .
	\]

	\textbf{Lagrange 乘数法（同一结果）.}
	令 $L=8 x y z-\lambda\left(\frac{x^2}{a^2}+\frac{y^2}{b^2}+\frac{z^2}{c^2}-1\right)$，由
	\[
	8 y z=\frac{2 \lambda x}{a^2},\qquad 8 x z=\frac{2 \lambda y}{b^2},\qquad 8 x y=\frac{2 \lambda z}{c^2}
	\]
	分别乘 $x$、$y$、$z$ 得 $8 x y z=\frac{2 \lambda x^2}{a^2}=\frac{2 \lambda y^2}{b^2}=\frac{2 \lambda z^2}{c^2}$，
	故 $\frac{x^2}{a^2}=\frac{y^2}{b^2}=\frac{z^2}{c^2}=\frac{1}{3}$，与上同。

	\par\vspace{1mm}\noindent
	\textbf{注：}{\kaishu
	几何意义：$8 x y z$ 正是内接于椭球的、各面平行于坐标面的长方体体积，
	其最大值为 $\frac{8 a b c}{3 \sqrt{3}}$，恰为椭球体积 $\frac{4}{3} \pi a b c$ 的 $\frac{2}{\sqrt{3} \pi}$ 倍。
	}\par
}


\item
设 $\ln \left(x^2+y\right)=x^3 y+\sin x$ 在 $x=0$ 处一阶可导，
求 $\left.\frac{d y}{d x}\right|_{x=0}$ 的值。

\banswer{
	\textbf{定初值.} 在方程中令 $x=0$ 得 $\ln y=0$，故 $y(0)=1$，对应点为 $(0,1)$，
	此处 $x^2+y=1>0$，$\ln$ 有意义。

	\textbf{两端对 $x$ 求导.}
	\[
	\frac{2 x+y^{\prime}}{x^2+y}=3 x^2 y+x^3 y^{\prime}+\cos x .
	\]

	\textbf{代入 $x=0$，$y=1$.}
	\[
	\frac{0+y^{\prime}(0)}{0+1}=0+0+1 \Longrightarrow \left.\frac{d y}{d x}\right|_{x=0}=1 .
	\]
}


\item
设函数 $f(x)$ 一阶可导，其反函数为 $g(x)$，且
$\int_0^{f(x)} g(t) d t=x^2 e^x$，$f(0)=0$，求 $f(x)$。

\banswer{
	记 $G(u)=\int_0^u g(t) d t$，则条件即 $G(f(x))=x^2 e^x$。两端对 $x$ 求导：
	\[
	g(f(x)) \cdot f^{\prime}(x)=\left(2 x+x^2\right) e^x .
	\]
	因 $g$ 是 $f$ 的反函数，$g(f(x))=x$，故
	\[
	x f^{\prime}(x)=x(2+x) e^x \Longrightarrow f^{\prime}(x)=(x+2) e^x \quad(x \neq 0),
	\]
	先在 $x>0$、$x<0$ 的相应区间分别积分（积分常数暂记为 $C_+$、$C_-$），得
	\[
	f(x)=\int(x+2) e^x d x=(x-1) e^x+2 e^x+C=(x+1) e^x+C .
	\]
	由 $f$ 在 $0$ 处连续且 $f(0)=0$，两侧积分常数均为 $-1$。所得显式函数在 $0$ 处也可导，且 $f^{\prime}(0)=2$。所以
	\[
	f(x)=(x+1) e^x-1 .
	\]

	\textbf{验证.} $f(0)=0$，$f^{\prime}(x)=(x+2) e^x$，
	于是 $\frac{d}{d x} G(f(x))=x \cdot(x+2) e^x=\left(x^2+2 x\right) e^x=\frac{d}{d x}\left(x^2 e^x\right)$，
	且两端在 $x=0$ 处同为 $0$，故在包含 $0$、使 $f$ 可逆的区间上 $G(f(x))=x^2 e^x$ 成立。例如可取 $x>-2$；该函数不能在整条实轴上定义单值反函数。
}


\item
求直线 $L:\left\{\begin{array}{l}x-y+z=1 \\ 2 x+y-z=1\end{array}\right.$
在平面 $\pi: x-2 y-2 z=1$ 上的投影直线方程 $L_0$，
以及投影直线 $L_0$ 绕 $z$ 轴旋转所得曲面的方程。

\banswer{
	\textbf{一、化 $L$ 为参数式.} 两方程相加得 $3 x=2$，即 $x=\frac{2}{3}$；
	代回第一式得 $z-y=\frac{1}{3}$。故
	\[
	L:\ \left(\frac{2}{3},\ t,\ t+\frac{1}{3}\right),\qquad \vec{s}=(0,1,1),
	\]
	取 $L$ 上一点 $P_0\left(\frac{2}{3}, 0, \frac{1}{3}\right)$。

	\textbf{二、求投影直线.} 作过 $L$ 且垂直于 $\pi$ 的平面 $\pi^*$，其法向量为
	\[
	\vec{n}^*=\vec{s} \times \vec{n}_\pi=(0,1,1) \times(1,-2,-2)=(0,1,-1).
	\]
	$\pi^*$ 过 $P_0$，方程为 $y-\left(z-\frac{1}{3}\right)=0$，即 $3 y-3 z+1=0$。于是
	\[
	L_0:\ \left\{\begin{array}{l}x-2 y-2 z=1 \\ 3 y-3 z+1=0\end{array}\right.
	\]
	化为参数式：由第二式 $z=y+\frac{1}{3}$，代入第一式得 $x=4 y+\frac{5}{3}$，故
	\[
	L_0:\ \left(4 t+\frac{5}{3},\ t,\ t+\frac{1}{3}\right),\qquad
	\text{即}\quad \frac{x-\frac{5}{3}}{4}=\frac{y}{1}=\frac{z-\frac{1}{3}}{1} .
	\]
	（校验：$P_0$ 在 $\pi$ 上的投影为
	$P_0-\frac{\vec{n}_\pi \cdot P_0-1}{\left|\vec{n}_\pi\right|^2} \vec{n}_\pi
	=\left(\frac{7}{9},-\frac{2}{9}, \frac{1}{9}\right)$，
	确在 $L_0$ 上（$t=-\frac{2}{9}$）；$\vec{s}$ 在 $\pi$ 上的投影 $\propto(4,1,1)$，与 $L_0$ 方向一致。）

	\textbf{三、绕 $z$ 轴旋转.} 在 $L_0$ 上用 $z$ 作参数：$t=z-\frac{1}{3}$，则
	\[
	x=4 z+\frac{1}{3},\qquad y=z-\frac{1}{3} .
	\]
	绕 $z$ 轴旋转时，$z$ 不变而 $x^2+y^2$ 保持为该高度处母线点到轴的距离平方，故
	\[
	x^2+y^2=\left(4 z+\frac{1}{3}\right)^2+\left(z-\frac{1}{3}\right)^2
	=17 z^2+2 z+\frac{2}{9},
	\]
	即旋转曲面为
	\[
	9 x^2+9 y^2-153 z^2-18 z-2=0 .
	\]
}


\item
计算曲面积分
\[
I=\iint\limits_{\Sigma} x^2 d y d z+y^2 d z d x+z^2 d x d y,
\]
其中 $\Sigma$ 为锥面 $x^2+y^2=z^2$ 上 $z \in(0,2)$ 的部分，取外侧。

\banswer{
	\textbf{参数化.} 令 $x=r \cos \theta$，$y=r \sin \theta$，$z=r$，
	$r \in[0,2]$，$\theta \in[0,2 \pi)$。由
	\[
	\vec{r}_r \times \vec{r}_\theta=(\cos \theta, \sin \theta, 1) \times(-r \sin \theta, r \cos \theta, 0)
	=(-r \cos \theta,\ -r \sin \theta,\ r)
	\]
	知它指向锥面内侧（远离 $z$ 轴的分量为负），故\textbf{外侧}对应
	\[
	d \vec{S}=(r \cos \theta,\ r \sin \theta,\ -r)\, d r\, d \theta .
	\]

	\textbf{计算.}
	\[
	I=\int_0^{2 \pi} \!\! \int_0^2\left[r^2 \cos ^2 \theta \cdot r \cos \theta
	+r^2 \sin ^2 \theta \cdot r \sin \theta+r^2 \cdot(-r)\right] d r\, d \theta .
	\]
	因 $\int_0^{2 \pi} \cos ^3 \theta\, d \theta=\int_0^{2 \pi} \sin ^3 \theta\, d \theta=0$，前两项为零，于是
	\[
	I=-\int_0^{2 \pi} d \theta \int_0^2 r^3 d r=-2 \pi \cdot 4=-8 \pi .
	\]

	\textbf{用 Gauss 公式校验.} 补上顶盖 $\Sigma_1: z=2$，$x^2+y^2 \leqslant 4$（取上侧），
	则 $\Sigma+\Sigma_1$ 围成锥体 $\Omega:\ x^2+y^2 \leqslant z^2$，$0 \leqslant z \leqslant 2$。
	$\Sigma_1$ 上 $d y d z=d z d x=0$，故
	$\iint_{\Sigma_1}=\iint_{x^2+y^2 \leqslant 4} 4\, d x d y=16 \pi$。而
	\[
	\iiint_{\Omega}(2 x+2 y+2 z) d V=0+0+2 \int_0^2 z \cdot \pi z^2 d z=8 \pi
	\]
	（前两项由对称性为零）。故 $I=8 \pi-16 \pi=-8 \pi$，与上一致。

	\par\vspace{1mm}\noindent
	\textbf{注：}{\kaishu
	回忆版题面中间一项写作 $y^2 d x d z$。由于 $\int_0^{2 \pi} \sin ^3 \theta\, d \theta=0$，
	不论取 $y^2 d z d x$ 还是 $y^2 d x d z=-y^2 d z d x$，该项贡献都为零，
	所以答案同为 $-8 \pi$。
	}\par
}


\end{enumerate}
